\documentclass{article}
\usepackage{amsmath,amssymb,amsthm}
\usepackage{algorithm,algorithmic}
\usepackage{hyperref}
\usepackage{enumitem}
\newtheorem{theorem}{Theorem}
\newtheorem{lemma}{Lemma}
\newtheorem{definition}{Definition}
\newtheorem{assumption}{Assumption}
\newcommand{\argmin}{\operatorname*{arg\,min}}
\newcommand{\diam}{\operatorname{diam}}
\begin{document}

\section*{Frank–Wolfe (Conditional Gradient) Algorithm}

\section*{Mathematical Formulation}
Let $f:\mathbb{R}^d\rightarrow\mathbb{R}$ be a convex and $L$‑smooth function and let $\mathcal{C}\subset\mathbb{R}^d$ be a non‑empty compact convex set.  
At iteration $k\ge 0$ the algorithm performs  

\[
\begin{aligned}
\text{(Linear minimisation oracle)}\qquad 
s_k &:= \argmin_{s\in\mathcal{C}}\;\langle \nabla f(x_k),\,s\rangle, \tag{1}\\[4pt]
\text{(Update)}\qquad 
x_{k+1} &:= (1-\gamma_k)\,x_k + \gamma_k\, s_k, \tag{2}
\end{aligned}
\]
where $\gamma_k\in[0,1]$ is a stepsize (either pre‑specified or obtained by a line‑search).

\section*{Pseudocode}
\begin{algorithm}[H]
\caption{Frank–Wolfe}
\begin{algorithmic}[1]
\STATE \textbf{Input:} convex $L$‑smooth $f$, feasible set $\mathcal{C}$, initial point $x_0\in\mathcal{C}$, stepsize rule $\{\gamma_k\}_{k\ge0}$, max iterations $K$.
\FOR{$k=0,1,\dots,K-1$}
    \STATE Compute gradient $g_k=\nabla f(x_k)$.
    \STATE \textbf{LMO:} $s_k \leftarrow \argmin_{s\in\mathcal{C}}\langle g_k,s\rangle$.
    \STATE Choose stepsize $\gamma_k\in[0,1]$ (e.g. $\gamma_k=2/(k+2)$ or exact line‑search).
    \STATE Update $x_{k+1}\leftarrow (1-\gamma_k)x_k+\gamma_k s_k$.
\ENDFOR
\STATE \textbf{Return} $x_K$ (or the best iterate w.r.t. the duality gap).
\end{algorithmic}
\end{algorithm}

\section*{Dry Run Example}
Consider the one‑dimensional problem  

\[
f(w)=(w-3)^2,\qquad \mathcal{C}=[0,5],\qquad w_0=0,\qquad \gamma_k\equiv \gamma=0.1 .
\]

\begin{enumerate}[label=Iteration \arabic*:, leftmargin=*]
\item \textbf{Iteration 0}
\[
\begin{aligned}
g_0 &= \nabla f(w_0)=2(w_0-3)=-6,\\
s_0 &= \argmin_{s\in[0,5]}\langle -6,s\rangle =\argmin_{s\in[0,5]}(-6s)=5,\\
w_1 &= (1-0.1) w_0 +0.1 s_0 =0.5,\\
f(w_1) &= (0.5-3)^2 = 6.25 .
\end{aligned}
\]

\item \textbf{Iteration 1}
\[
\begin{aligned}
g_1 &= 2(w_1-3)=2(0.5-3)=-5,\\
s_1 &= \argmin_{s\in[0,5]}(-5s)=5,\\
w_2 &= 0.9\cdot 0.5 +0.1\cdot5 =0.95,\\
f(w_2) &= (0.95-3)^2 = 4.2025 .
\end{aligned}
\]

\item \textbf{Iteration 2}
\[
\begin{aligned}
g_2 &= 2(w_2-3)=2(0.95-3)=-4.1,\\
s_2 &= \argmin_{s\in[0,5]}(-4.1 s)=5,\\
w_3 &= 0.9\cdot0.95+0.1\cdot5 =1.355,\\
f(w_3) &= (1.355-3)^2 \approx 2.701 .
\end{aligned}
\]
\end{enumerate}
The iterates move monotonically towards the optimum $w^\star=3$ while staying inside $\mathcal{C}$.

\newpage
\section*{Assumptions}
\begin{assumption}[Convexity and Smoothness]\label{ass:conv-smooth}
$f$ is convex and differentiable on an open set containing $\mathcal{C}$, and its gradient is $L$‑Lipschitz:
\[
\|\nabla f(x)-\nabla f(y)\|\le L\|x-y\|\quad\forall x,y\in\mathcal{C}.
\]
\end{assumption}
\begin{assumption}[Compact Feasible Set]\label{ass:compact}
$\mathcal{C}$ is compact and convex with finite diameter 
\[
D:=\diam(\mathcal{C}) = \max_{x,y\in\mathcal{C}}\|x-y\| <\infty .
\]
\end{assumption}
\begin{assumption}[Linear Minimisation Oracle]\label{ass:lmo}
For any $g\in\mathbb{R}^d$, the oracle returns an exact solution $s\in\mathcal{C}$ of (1). 
\end{assumption}

\section*{Main Convergence Theorem}
\begin{theorem}[Primal Gap Decay]\label{thm:fw-rate}
Let $\{x_k\}$ be generated by the Frank–Wolfe algorithm with stepsizes $\gamma_k= \frac{2}{k+2}$. Under Assumptions \ref{ass:conv-smooth}–\ref{ass:lmo} the primal suboptimality satisfies
\[
\boxed{ \; f(x_k)-f(x^\star) \le \frac{2 L D^{2}}{k+2}\; } \qquad\forall k\ge 0,
\]
where $x^\star\in\arg\min_{x\in\mathcal{C}} f(x)$.
\end{theorem}

\section*{Theoretical Properties}
\begin{itemize}
\item \textbf{Stability.} The iterates remain in $\mathcal{C}$ for any $\gamma_k\in[0,1]$ because $\mathcal{C}$ is convex.
\item \textbf{Stepsize Sensitivity.} The standard schedule $\gamma_k=2/(k+2)$ yields the $O(1/k)$ rate. Using a constant $\gamma\in(0,1)$ still guarantees convergence but with a slower linear rate that depends on the curvature constant (see Lemma~\ref{lem:curv-bound}).
\item \textbf{Comparison.} Compared with projected gradient descent, Frank–Wolfe avoids costly Euclidean projections and attains the same $O(1/k)$ rate for smooth convex problems, at the expense of a potentially larger constant $LD^2$.
\end{itemize}

\section*{Discussion}
The bound in Theorem~\ref{thm:fw-rate} is tight for worst‑case smooth convex functions; improvements are possible under stronger curvature assumptions (e.g., strong convexity) or by using away‑steps/fully‑corrective variants.

\newpage
\section*{Preliminaries (Definitions and Lemmas)}
\begin{definition}[Curvature Constant]\label{def:curv}
The curvature of $f$ on $\mathcal{C}$ is
\[
C_f := \sup_{\substack{x,s\in\mathcal{C}\\ \gamma\in[0,1]}}
\frac{2}{\gamma^{2}}\Bigl[f\bigl(x+\gamma(s-x)\bigr)-f(x)-\gamma\langle\nabla f(x),s-x\rangle\Bigr].
\]
\end{definition}
\begin{lemma}[Smoothness Upper Bound]\label{lem:smooth}
If $f$ is $L$‑smooth then for any $x,s\in\mathcal{C}$ and $\gamma\in[0,1]$,
\[
f\bigl(x+\gamma(s-x)\bigr)\le f(x)+\gamma\langle\nabla f(x),s-x\rangle+\frac{L\gamma^{2}}{2}\|s-x\|^{2}. \tag{3}
\]
\end{lemma}
\begin{proof}
By $L$‑smoothness,
\[
f(y)\le f(x)+\langle\nabla f(x),y-x\rangle+\frac{L}{2}\|y-x\|^{2}\quad\forall x,y.
\]
Set $y=x+\gamma(s-x)$ to obtain (3). \hfill $\square$
\end{proof}
\begin{lemma}[Bound on Curvature]\label{lem:curv-bound}
For an $L$‑smooth function on a set of diameter $D$, the curvature satisfies $C_f\le L D^{2}$.
\end{lemma}
\begin{proof}
From (3) we have for any $x,s$ and $\gamma$,
\[
\frac{2}{\gamma^{2}}\Bigl[f(x+\gamma(s-x))-f(x)-\gamma\langle\nabla f(x),s-x\rangle\Bigr]
\le L\|s-x\|^{2}\le L D^{2}.
\]
Taking the supremum over $x,s,\gamma$ yields $C_f\le L D^{2}$. \hfill $\square$
\end{proof}

\section*{Main Proof (Step‑by‑Step Derivation)}
We prove Theorem~\ref{thm:fw-rate} using the curvature constant.

\noindent\textbf{Notation.} Let $x^\star\in\arg\min_{\mathcal{C}}f$ and define the \emph{Frank–Wolfe gap}
\[
g(x_k):=\max_{s\in\mathcal{C}}\langle \nabla f(x_k),x_k-s\rangle
      =\langle \nabla f(x_k),x_k-s_k\rangle .
\]

\begin{proof}
\begin{enumerate}
\item[Step 1.] Apply Lemma~\ref{lem:smooth} with $x=x_k$, $s=s_k$, and $\gamma=\gamma_k$:
\[
f(x_{k+1})\le f(x_k)+\gamma_k\langle\nabla f(x_k),s_k-x_k\rangle
                +\frac{L\gamma_k^{2}}{2}\|s_k-x_k\|^{2}. \tag{4}
\]

\item[Step 2.] By definition of $s_k$,
\[
\langle\nabla f(x_k),s_k-x_k\rangle = - g(x_k). \tag{5}
\]

\item[Step 3.] Since $\mathcal{C}$ has diameter $D$, $\|s_k-x_k\|\le D$. Insert (5) and the bound $\|s_k-x_k\|^{2}\le D^{2}$ into (4):
\[
f(x_{k+1})\le f(x_k)-\gamma_k g(x_k)+\frac{L D^{2}}{2}\gamma_k^{2}. \tag{6}
\]

\item[Step 4.] Use the convexity of $f$:
\[
f(x_k)-f(x^\star)\le \langle\nabla f(x_k),x_k-x^\star\rangle
                 \le g(x_k), \tag{7}
\]
because $x^\star\in\mathcal{C}$ and $g(x_k)$ is the maximal inner product over $\mathcal{C}$.

\item[Step 5.] Combine (6) and (7):
\[
f(x_{k+1})-f(x^\star)
   \le f(x_k)-f(x^\star)-\gamma_k\bigl[f(x_k)-f(x^\star)\bigr]
        +\frac{L D^{2}}{2}\gamma_k^{2}. \tag{8}
\]

\item[Step 6.] Rearrange (8) to obtain a recursion for the primal gap $\Delta_k:=f(x_k)-f(x^\star)$:
\[
\Delta_{k+1}\le (1-\gamma_k)\Delta_k+\frac{L D^{2}}{2}\gamma_k^{2}. \tag{9}
\]

\item[Step 7.] Choose the standard diminishing stepsize $\displaystyle\gamma_k=\frac{2}{k+2}$.
Plugging into (9) yields
\[
\Delta_{k+1}\le\left(1-\frac{2}{k+2}\right)\Delta_k
               +\frac{L D^{2}}{2}\left(\frac{2}{k+2}\right)^{2}
            =\frac{k}{k+2}\Delta_k+\frac{2 L D^{2}}{(k+2)^{2}}. \tag{10}
\]

\item[Step 8.] Multiply both sides of (10) by $(k+1)(k+2)$:
\[
(k+1)(k+2)\Delta_{k+1}
   \le k(k+1)\Delta_k+2 L D^{2}. \tag{11}
\]

\item[Step 9.] Define $A_k:=(k+1)\Delta_k$.  Then (11) becomes
\[
A_{k+1}\le A_k+\frac{2 L D^{2}}{k+2}. \tag{12}
\]

\item[Step 10.] Unfold the recursion from $0$ to $k$:
\[
A_{k+1}\le A_0+2 L D^{2}\sum_{t=0}^{k}\frac{1}{t+2}
        \le A_0+2 L D^{2}\bigl(1+\ln(k+2)\bigr). \tag{13}
\]
Since $A_0=(0+1)\Delta_0\le f(x_0)-f(x^\star)$ is finite, the dominant term for large $k$ is the harmonic sum.

\item[Step 11.] A tighter bound is obtained by observing that the recursion (9) is a linear contraction; solving it directly gives
\[
\Delta_{k}\le \frac{2 L D^{2}}{k+2}. \tag{14}
\]
Indeed, we prove (14) by induction.  
Base $k=0$: $\Delta_0\le \frac{2LD^{2}}{2}=LD^{2}$ holds because $\Delta_0\le \frac{L}{2}D^{2}$ from smoothness.  
Inductive step: assume $\Delta_k\le \frac{2LD^{2}}{k+2}$. Using (9) with $\gamma_k=2/(k+2)$,
\[
\Delta_{k+1}\le\left(1-\frac{2}{k+2}\right)\frac{2LD^{2}}{k+2}
               +\frac{LD^{2}}{2}\left(\frac{2}{k+2}\right)^{2}
            =\frac{2LD^{2}}{k+3},
\]
which completes the induction. \hfill $\square$ \qedhere
\end{enumerate}
\end{proof}

\section*{Rate Derivation (With Constants)}
From (14) we have the explicit $O(1/k)$ rate:
\[
f(x_k)-f(x^\star)\le \frac{2 L D^{2}}{k+2}\qquad\forall k\ge0.
\]
The constant $2LD^{2}$ depends only on the smoothness $L$ and the diameter $D$ of $\mathcal{C}$.

\section*{Special Cases}
\begin{enumerate}
\item \textbf{Strongly Convex $f$.} If $f$ is $\mu$‑strongly convex, the primal gap contracts linearly:
\[
\Delta_{k}\le\left(1-\frac{\mu\gamma_k}{L}\right)^{k}\Delta_0,
\]
when a constant stepsize $\gamma_k\equiv\gamma\in(0,1]$ is used.
\item \textbf{Exact Line‑Search.} Choosing $\gamma_k=\arg\min_{\gamma\in[0,1]}f\bigl(x_k+\gamma(s_k-x_k)\bigr)$ yields the same $O(1/k)$ bound but with a possibly smaller empirical constant.
\item \textbf{Away‑Step Variant.} Incorporating away steps reduces the dependence on $D$ to the \emph{pyramidal width} of $\mathcal{C}$ and can achieve linear convergence for polytopes.
\end{enumerate}

\end{document}